% Introduction to the EU - Week 9 Reading Summary
% Compile: pdflatex Some-remarks-rule-of-law-Poland-Hungary-reading-summary.tex

\documentclass[11pt,a4paper]{article}
\usepackage[utf8]{inputenc}
\usepackage[T1]{fontenc}
\usepackage[margin=2.5cm]{geometry}
\usepackage{booktabs}
\usepackage{enumitem}
\setlist{nosep,leftmargin=*}
\usepackage{parskip}
\usepackage{microtype}
\usepackage{hyperref}
\usepackage{xcolor}
\definecolor{eublue}{RGB}{0,51,153}

\hypersetup{
  colorlinks=true,
  linkcolor=eublue,
  urlcolor=eublue,
  pdftitle={Rule of Law in Poland and Hungary Reading Summary},
  pdfauthor={Introduction to the EU, UCM}
}

\begin{document}

\title{Walter Rech: \textit{Some remarks on the EU's action on the erosion of the rule of law in Poland and Hungary}\\Reading Summary}
\author{Introduction to the European Union\\Universidad Complutense de Madrid\\Week 9}
\date{}
\maketitle
\vspace{-1em}
\hrule
\vspace{1em}

\section*{Bibliographic Information}
\begin{itemize}
\item \textbf{Author:} Walter Rech
\item \textbf{Article:} Some remarks on the EU's action on the erosion of the rule of law in Poland and Hungary
\item \textbf{Journal:} Journal of Contemporary European Studies
\item \textbf{Year/Volume/Pages:} 2018, 26(3), 334--345
\item \textbf{DOI:} \url{https://doi.org/10.1080/14782804.2018.1498770}
\item \textbf{Source file:} Some remarks on the EU s action on the erosion of the rule of law in Poland and Hungary.pdf
\end{itemize}

\section*{Thesis}
Rech supports criticism of democratic backsliding in Hungary and Poland, but argues that EU rule-of-law discourse is internally ambivalent and inconsistently enforced; without self-critical reflection, liberal-democratic rhetoric becomes politically vulnerable and less effective.

\section*{Structure Of The Argument}
\subsection*{I. Foundational Or Subordinate?}
\begin{itemize}
\item EU formally treats rule of law/democracy as foundational (Art. 2 TEU, Copenhagen criteria, Rule of Law Framework).
\item In practice, enforcement is uneven and frequently subordinated to geopolitical/economic priorities.
\item Article 7 sanctions are politically hard to execute due to unanimity constraints and bloc shielding dynamics.
\end{itemize}

\subsection*{II. Why Fidesz And PiS Stay Politically Effective}
\begin{itemize}
\item Backsliding is tied to institutional redesign, not only populist rhetoric.
\item Domestic support is reinforced through material-economic policies and identity politics.
\item EU pressure can be reframed as external hostility, strengthening nationalist narratives.
\end{itemize}

\subsection*{III. Theoretical Ambivalence Of Rule-Of-Law Language}
\begin{itemize}
\item Rule of law and democracy are not fixed, universally interpreted categories in practice.
\item EU narratives can become self-righteous if they ignore internal contradictions.
\item Liberal-democratic legitimacy should remain critical and reflexive rather than essentialized.
\end{itemize}

\section*{Policy And Strategy Implications}
\begin{itemize}
\item \textbf{Sanctions dilemma:} punitive responses may polarize and hurt citizens more than elites.
\item \textbf{Economic leverage:} conditionality can be effective but raises legitimacy/fairness concerns.
\item \textbf{Institutional challenge:} EU needs credible, consistent enforcement plus political vision beyond legal formalism.
\end{itemize}

\section*{Relevance To Course Themes}
\begin{itemize}
\item Clarifies the gap between constitutional ideals and practical governance in the EU.
\item Links legal integration debates to domestic political economy and sovereignty claims.
\item Helps evaluate whether EU authority is legal, political, or both in rule-of-law crises.
\end{itemize}

\section*{Discussion Questions}
\begin{enumerate}
\item Can the EU enforce rule-of-law standards consistently without appearing coercive or selective?
\item Is conditionality of EU funds a legitimate democratic tool or a socially regressive sanction?
\item How can EU institutions defend liberal democracy while acknowledging internal contradictions?
\end{enumerate}

\vspace{1em}
\noindent\footnotesize\textit{Introduction to the European Union, Universidad Complutense de Madrid.}

\end{document}
